\section{Leakage boundaries and exact bounded-decoder conventions}\label{sec:supp-r70}
The transient-core lemma in the main article uses a nonzero response tending to zero, a prescribed initial label and a decoder in $[0,1]$. The assumptions ensure that the initial label belongs to the output-visible core even when its immediate output is zero. The positive path argument tests every active entry, not only the norm of one arbitrarily chosen physical state. The appended zero-output label remains live until the productive herald. Removing it by declaring a hidden leakage event terminal would change the experiment.

\begin{proposition}[Degenerate leakage parameters]\label{prop:r70-degenerate}
For $u_n=\lambda^n[1+\gamma\alpha^n\cos(n\vartheta+\varphi)]/2$, interpret the initial term separately at zero parameters. With $0<\lambda<1$, zero contrast gives least live memory two. At $\alpha=0$ the least size is two if $\gamma\cos\varphi=0$, and three otherwise. At $\lambda=0$ it is one when $u_0=0$ and two otherwise. At $\lambda=1$ there is no mandatory leakage label: the inherited unattenuated theorem applies in its stated nonreal-mode case. Background $\eta=1$ gives the constant response $1/2$ and needs one label.
\end{proposition}
\begin{proof}
For zero contrast, the response is $\lambda^n/2$. One live label cannot change its output; a transient label of retention $\lambda$ and output $1/2$, followed by an absorbing zero-output label, realizes it with two.

At $\alpha=0$, let $a=u_0=(1+\gamma\cos\varphi)/2$. The subsequent terms are $\lambda^n/2$, $n\ge1$. A three-label construction has a starting label of output $a$, leading deterministically on its first failure to a transient label with retention $\lambda$ and output $\lambda/2$, followed by a zero-output label. Its response is exactly the required one. A two-label realization has, by Lemma~\ref{lem:r70-core}, a one-dimensional active core, so its entire response must be $a r^n$. Matching the positive terms at $n=1,2$ forces $r=\lambda$ and $a=1/2$. If $a=0$, nonzero later output already excludes a one-dimensional core. Thus two are possible precisely in the stated case.

At $\lambda=0$ only the initial output can be nonzero. A starting label of that output followed by a zero-output label suffices; one label is possible precisely when all terms are zero. At $\lambda=1$ the output no longer has the decay hypothesis used to force a zero-output label. The remaining claims follow directly from the constant response or the inherited theorem.
\end{proof}

The theorem's nonreal-mode assumption also matters. If $\vartheta$ is a multiple of $\pi$, exponential modes can coincide or cancel. One must minimize the resulting actual response rather than apply a nonreal spectral count. No claim about those coincidences is hidden in the four-mode Hankel statement.

\subsection{The physical audit is unchanged by the realization argument}
For the six-dimensional squared embedding, the herald probability is $p$ in every model; it is not the leakage probability. Its first compulsory call and every failed call cost one. Hence the least cost is $1/p$ and every boundary-feasible policy accepts the first herald. The flag parameter $t$ changes an unobserved flag channel, not the herald hazard. The two signs are physical audit-model parameters with the same fixed loss $(r-z)^2$, not a choice of loss after observing the answer. The three effects in the inherited construction have positive mass and sum to the identity. Their computed moments give
\[
 \mathbb E_{t,\sigma}(r-Z)^2
 =2+2\varepsilon\sigma p\sum_{n\ge0}(qt)^n(b_n-u_n).
\]
Averaging the two signs gives two for every rule. Thus the target response attains two even on the finite witness family, while exact finite-memory equality requires both opposed inequalities at every test parameter. The hidden sink is never supplied as an observer report.

\subsection{What the finite test certifies}
The generating function of an $M$-label stochastic decoder has denominator $\det(I-zQ)$ and numerator degree at most $M-1$. The determinant is nonzero throughout $[0,q\bar t]$ because $q\bar t<1$. The target's denominator is likewise nonzero there. Consequently the $M+r$-point theorem is a polynomial identity argument with a known degree bound. It does not certify an arbitrary fitted function from finitely many samples, nor identify an unrestricted full-history rule from those finitely many equations.

For the leaking irrational-phase example, the exact fixed-$M$ optimization can be encoded with variables $Q,d,v_j,w$ and the inequalities
\[
 Q\ge0,\ Q\one=\one,\quad 0\le d\le\one,\quad
 (I-qt_jQ)v_j=d,\quad
 -w\le(v_j)_1-U(qt_j)\le w.
\]
The minimizing $w$ exists and is strictly positive. Quantifier elimination over this compact graph computes it as an algebraic number and returns a minimizing controller. The proof supplies this terminating procedure for every prescribed size; the finite diagnostics do not claim to have run a global elimination for all sizes. The number of model pairs grows with $M$, which is essential to the scope of the finite witness assertion.

\begin{proposition}[An exact finite-witness one-label value]\label{prop:r70-one-label}
In Theorem~\ref{thm:r70-finite-tests}, take $M=1$, the four parameters $t=0,1/6,1/3,1/2$, and both audit signs. The least one-live-label risk is
\[
 2+\frac{331}{13573},
\]
attained by the constant probability $d=12249/13573$ of decoder $+1$.
\end{proposition}
\begin{proof}
The generating target is
\[
 U(z)=\frac{-8z^2+28z-40}{5z^3-22z^2+44z-40},
 \qquad U(0)=1,\quad U(1/4)=2144/1939.
\]
Every one-label rule is a constant $d\in[0,1]$. Its uniform resolvent error is at least
$\max\{1-d,\,4d/3-2144/1939\}$ from the two endpoint tests. Their intersection gives $d_*=12249/13573$ and lower error $1324/13573$. Direct rational substitution at $z=1/12$ and $z=1/6$ gives errors respectively
$-465602268/9412509029$ and $6704202/487474295$, whose absolute values are strictly below this lower error. Thus all four tests are bounded by the endpoint value at $d_*$. Finally $2\varepsilon p=1/4$ converts resolvent error to the stated excess squared risk. This proves optimality over all stochastic one-label decoders; no local numerical optimizer is used.
\end{proof}

\subsection{Background, initial mixing and precision}
The background effect $E_\eta=\eta I/2+(1-\eta)E$ is a valid actual effect on the qutrit, with eigenvalues in $[\eta/2,1-\eta/2]$. Its positive floor changes which terminal responses must vanish on the stationary face. It does not add a report, a reset, an observed physical state or a free random seed.

The construction stores the $T$ exceptional prefix values. Only its final failure row chooses one of two disjoint suffix machines, and that choice is subsequently represented by the retained label. Their average realizes the complete suffix from its first step. This is why the stated count is $T+2+m(\lambda\alpha,\vartheta)$, with no additional initialization label and no uncounted intercall mixture memory.

For algebraic input the background target has an algebraic linear representation of dimension at most four. A rationally certified polygon or the existing spectral size theorem supplies a finite size upper. The prescribed-size algebraic equality test then terminates a minimal-size search. Exact algebraic encodings and their sampling cost are separate from the live-label count. The upper logarithmic dependence on $1/\eta$ is not a matched precision or program-bit lower bound.
